\documentclass[11pt]{article}

% --- Packages ---
\usepackage[margin=1in]{geometry}
\usepackage{setspace}
\onehalfspacing
\usepackage{amsmath,amssymb}
\usepackage{booktabs}
\usepackage{graphicx}
\usepackage[hidelinks]{hyperref}
\usepackage{natbib}
\usepackage{caption}
\usepackage{enumitem}
\usepackage{float}
\usepackage{url}

% --- Title ---
\title{Capacity Portfolio Evaluation Under Multi-Dimensional Uncertainty:\\A Stochastic Simulation Study}
\author{Omar Mourssi --- 1008311961}
\date{April 15, 2026}

\begin{document}
\maketitle

%==========================================================================
% SECTION 1: INTRODUCTION
%==========================================================================
\section{Introduction}\label{sec:intro}

Long-term electricity capacity planning requires selecting a resource mix that balances cost and reliability across plausible future conditions. In Ontario, the IESO projects a 75\% increase in provincial electricity demand by 2050, driven by electrification, electric-vehicle supply chains, data centres, and industrial expansion \citep{IESO2025}. As the generation mix shifts toward variable renewables and battery storage, deterministic planning methods that evaluate a small number of fixed scenarios may not reveal when uncertainty changes which resource mix performs best.

Recent literature emphasizes explicit uncertainty treatment in capacity planning. \citet{Roald2023} review optimization under uncertainty in power systems. \citet{Stover2023} argue that deterministic reserve-margin rules are insufficient for high-renewable systems. \citet{Wang2024} show that temporal correlation between demand, renewables, and storage materially affects capacity credit. \citet{Leibowicz2024} find that capturing operational detail in Monte Carlo adequacy simulation changes which resources are identified as adequate.

A particularly relevant contribution is \citet{ParkBaldick2016}, who formulate a multi-year stochastic capacity expansion model and show that uncertainty in demand and wind availability changes optimal build timing and resource mix. The present study adapts this uncertainty-aware logic to a different problem: rather than optimizing expansion decisions, it evaluates a systematically generated set of fixed candidate portfolios using Monte Carlo simulation and formal statistical selection. Park and Baldick solve a multi-stage stochastic program; this project takes candidate mixes as given and asks which performs best under uncertainty, using Common Random Numbers (CRN), the Kim--Nelson (KN) sequential elimination procedure \citep{KimNelson2001}, and Bonferroni-corrected confidence intervals.

Candidate portfolios are generated by discretizing the feasible allocation simplex over firm generation, renewable generation, and battery storage at 10\% increments subject to a fixed total capacity budget. The simulation model is a stylized one-zone system, not calibrated to a specific jurisdiction, with input distributions calibrated from Ontario empirical data (IESO demand records and the CODERS dataset \citep{CODERS2025}).

The remainder of this report is organized as follows. Section~\ref{sec:problem} defines the problem and performance measure. Section~\ref{sec:model} describes the simulation model. Section~\ref{sec:design} presents the experimental design. Section~\ref{sec:results} reports the results. Section~\ref{sec:conclusion} discusses limitations and future work.

%==========================================================================
% SECTION 2: PROBLEM DESCRIPTION
%==========================================================================
\section{Problem Description}\label{sec:problem}

\subsection{Portfolio Generation}

The system comprises three technology classes: firm generation (nuclear, flexible firm, mid-merit gas, gas peakers), variable renewables (wind, solar), and battery storage. Total installed capacity is fixed at $C_{\text{total}} = 12{,}000$~MW, derived from 8,000~MW peak demand and a 50\% reserve margin. Each portfolio $P_i$ is defined by $(f_i^{\text{firm}}, f_i^{\text{renew}}, f_i^{\text{stor}})$ with $f_i^{\text{firm}} + f_i^{\text{renew}} + f_i^{\text{stor}} = 1$.

The feasible simplex is discretized at 10\% increments subject to $f^{\text{firm}} \in [0.20, 0.70]$, $f^{\text{renew}} \in [0.10, 0.60]$, $f^{\text{stor}} \in [0.10, 0.30]$, yielding $k = 16$ portfolios (Appendix~\ref{app:portfolios}). This systematic approach avoids selection bias and ensures uniform coverage of the design space.

\subsection{Performance Measure}

For portfolio $P_i$ in replication $j$:
\begin{equation}\label{eq:Y}
  Y_{ij} = F_i + C_{ij} + \lambda \, U_{ij},
\end{equation}
where $F_i$ is the annualized fixed cost (deterministic), $C_{ij}$ is simulated annual operating cost, $U_{ij}$ is unserved energy (MWh), and $\lambda = \$10{,}000$/MWh is the value of lost load. Operating cost is computed from hourly dispatch:
\begin{equation}\label{eq:C}
  C_{ij} = \sum_{t=1}^{8760} \sum_{g \in \mathcal{G}} c_g \, x_{g,i,j,t},
\end{equation}
where $c_g$ is marginal cost and $x_{g,i,j,t}$ is dispatched output in hour $t$. Both $C_{ij}$ and $U_{ij}$ emerge from the simulation described in Section~\ref{sec:model}; they are not sampled directly. The objective is $i^* = \arg\min_i \, \theta_i$ where $\theta_i = \mathbb{E}[Y_i]$.

\subsection{Research Questions}

\begin{enumerate}[leftmargin=*, itemsep=2pt]
  \item Which portfolio minimizes expected total cost under uncertainty, and does this differ from the deterministic recommendation?
  \item Which uncertainty factors materially affect portfolio ranking?
\end{enumerate}

%==========================================================================
% SECTION 3: MODEL DESCRIPTION
%==========================================================================
\section{Model Description}\label{sec:model}

\subsection{System Architecture}

Each portfolio allocates capacity across four firm sub-types, two renewable technologies, and battery storage. Firm generators are dispatched in merit order (Table~\ref{tab:merit}); renewables have zero marginal cost; storage follows a rule-based heuristic. Full parameter values are in Appendix~\ref{app:params}.

\begin{table}[h]
\centering
\caption{Firm generator merit-order stack.}\label{tab:merit}
\begin{tabular}{lcccc}
\toprule
Sub-type & Share & MC (\$/MWh) & FOR & Capital (\$/kW-yr) \\
\midrule
Nuclear       & 20\% & 22 & 5\% & 400 \\
Flexible firm & 35\% & 35 & 6\% & 100 \\
Mid-merit gas & 30\% & 50 & 5\% & 70 \\
Gas peaker    & 15\% & 80 & 8\% & 70 \\
\bottomrule
\end{tabular}
\end{table}

Nuclear marginal cost of \$22/MWh reflects the variable operating cost of refurbished Ontario nuclear units (fuel plus variable O\&M), not the all-in levelized cost of electricity.

\subsection{Hourly Dispatch Algorithm}

Each replication simulates 8,760 hours of system operation for a given portfolio. The dispatch model uses a sequential priority structure: renewable output is absorbed first (zero marginal cost), storage responds to net-load conditions, and firm generators are dispatched in ascending marginal-cost order to serve any remaining demand. This merit-order approach reflects standard practice in power system operations, where the cheapest available generation is dispatched first to minimize operating cost. The model does not co-optimize storage and firm dispatch; this simplification keeps the project tractable while preserving the core cost and adequacy dynamics that differentiate portfolios.

For each hour $t$, the dispatch proceeds as follows:

\begin{enumerate}[leftmargin=*, itemsep=2pt]
  \item \textbf{Demand.} $D(t) = D_{\text{base}}(t)(1 + \sigma_d \varepsilon(t))$, $\varepsilon(t) \sim N(0,1)$, $\sigma_d = 0.09$. The noise term is applied multiplicatively so that demand variability scales with load level; absolute deviations are larger during peak hours than off-peak. The iid assumption is a simplification noted in Section~\ref{sec:conclusion}.
  \item \textbf{Renewables.} Wind and solar capacity factors are drawn from NORTA AR(1) Beta distributions (Section~\ref{sec:inputs}). Aggregate renewable output is $R(t) = \text{MW}_w \cdot \text{CF}_w(t) + \text{MW}_s \cdot \text{CF}_s(t)$.
  \item \textbf{Residual demand.} $D_{\text{res}}(t) = \max(0,\; D(t) - R(t))$. This is the load that must be served by storage and firm generation.
  \item \textbf{Storage.} The battery discharges when residual demand exceeds the 70th percentile of the base demand profile, and charges from surplus generation during off-peak hours (residual demand below the 30th percentile). These thresholds are stylized operating rules chosen to approximate the economic intuition that storage should discharge during high-value periods and charge when energy is cheap. Discharge is limited by the minimum of current state of charge, rated power capacity, and remaining residual demand. Charging is limited by available surplus and rated power, subject to 85\% round-trip efficiency. A more sophisticated model would co-optimize storage dispatch across the full horizon; this simplification is acknowledged as a limitation.
  \item \textbf{Firm dispatch.} Available firm generators are dispatched in merit order (Table~\ref{tab:merit}). For each sub-type $g$, dispatched output is $x_g(t) = \min(\text{cap}_g \cdot a_g(t),\; D_{\text{rem}}(t))$, where $a_g(t)$ is the outage-adjusted availability fraction. The algorithm walks up the merit-order stack, accumulating operating cost as $x_g(t) \cdot c_g$ and subtracting dispatched output from remaining demand.
  \item \textbf{Unserved energy.} $u(t) = \max(0,\; D_{\text{rem}}(t))$ after all dispatch. Any demand that cannot be served after exhausting renewables, storage, and all available firm generation is recorded as unserved energy for that hour.
\end{enumerate}

\paragraph{Forced outage modeling.}
Each firm sub-type is represented by 10 discrete units of equal size. For each sub-type on each day, each unit independently draws a Bernoulli$(1 - \text{FOR})$ availability indicator, producing discrete availability steps of 0\%, 10\%, \ldots, 100\% of sub-type capacity. Outage states persist for the full 24-hour day. This per-unit, per-day granularity is coarser than a continuous-time Markov model but captures the daily availability risk that drives adequacy outcomes. The discrete 10-unit resolution limits the model's ability to represent partial outage states finely; this is noted as a limitation in Section~\ref{sec:conclusion}.

Annual outputs are accumulated over all 8,760 hours: $C_{ij} = \sum_t \sum_g c_g x_{g,t}$, \; $U_{ij} = \sum_t u(t)$, \; and \; $Y_{ij} = F_i + C_{ij} + \lambda U_{ij}$.

\subsection{Input Models and Empirical Calibration}\label{sec:inputs}

Input distributions are calibrated from IESO hourly Ontario demand \citep{IESO2025} and CODERS hourly capacity factors for 2,984 Ontario generation sites \citep{CODERS2025}.

\paragraph{Demand.}
The base profile $D_{\text{base}}(t)$ is fitted via OLS to 2024 IESO data using two seasonal harmonics (annual and semi-annual, capturing Ontario's winter/summer demand peaks) plus a diurnal harmonic, achieving $R^2 = 0.59$. The noise parameter $\sigma_d = 0.09$ is the standard deviation of relative residuals $(D_{\text{obs}} - D_{\text{base}})/D_{\text{base}}$.

\paragraph{Wind.}
Hourly aggregate CFs are fitted to Beta$(6.73, 15.69)$ via MLE (KS $= 0.021$, mean $= 0.300$). To preserve the strong temporal persistence (empirical lag-1 ACF $\hat\rho = 0.995$), CFs are generated via NORTA with AR(1) dependence:
\begin{equation}\label{eq:norta}
  Z(t) = \rho Z(t{-}1) + \sqrt{1-\rho^2}\,\epsilon(t), \quad
  \text{CF}(t) = F_{\text{Beta}}^{-1}(\Phi(Z(t))).
\end{equation}
This preserves the Beta marginal while inducing hourly autocorrelation. The distinction from iid sampling is consequential: iid draws overestimate hourly variability while underestimating sustained low-output periods that drive adequacy risk (Appendix~\ref{app:inputs}).

\paragraph{Solar.}
$\text{CF}_s(t) = E(t) \cdot K(t)$, where $E(t)$ is a sinusoidal diurnal envelope (hours 7--21) and $K(t) \sim \text{Beta}(1.77, 6.91)$ via NORTA AR(1) with $\rho = 0.928$, calibrated from CODERS.

\paragraph{Forced outages.}
Each firm sub-type has 10 units with independent daily Bernoulli$(1 - \text{FOR})$ draws, producing discrete availability steps of 0\%, 10\%, \ldots, 100\%.

%==========================================================================
% SECTION 4: EXPERIMENTAL DESIGN
%==========================================================================
\section{Experimental Design}\label{sec:design}

\subsection{Common Random Numbers}

All portfolios share identical random inputs within each replication, implemented via NumPy's \texttt{SeedSequence} spawning. Three independent streams (demand, outages, renewables) are deterministically generated from a base seed and replication index. CRN induces positive covariance between outputs, reducing $\text{Var}(Y_i - Y_\ell)$. Empirically, variance reduction ranges from 2$\times$ to 35$\times$ (Appendix, Figure~\ref{fig:crn}).

\subsection{Kim--Nelson Sequential Elimination}

The KN procedure \citep{KimNelson2001} selects the portfolio with smallest $\theta_i$, guaranteeing $P(\text{CS}) \geq 1 - \alpha$ when the best system is at least $\delta$ better than the second-best. Parameters: $\alpha = 0.05$, $\delta = \$50$M, $n_0 = 30$. The procedure initializes with $n_0$ CRN replications from all $k = 16$ portfolios, computes pairwise difference variances $S^2_{i\ell}$ (which automatically capture CRN covariance), then sequentially eliminates portfolios whose running means exceed a competitor's beyond a shrinking threshold. The procedure terminates when one survivor remains.

\subsection{Deterministic Benchmark Gap}

The deterministic recommendation $P_{\text{DET}}$ is obtained by setting all inputs to expected values and selecting the lowest-$Y$ portfolio. $P_{\text{DET}}$ and $P_{\text{KN}}$ are then evaluated for 200 paired CRN replications. The DBG $= \bar{Y}_{\text{DET}} - \bar{Y}_{\text{KN}}$ quantifies the cost of ignoring uncertainty in portfolio selection. This is not the classical VSS (which compares stochastic and deterministic optimizations); since we select from fixed portfolios, the DBG measures the value of uncertainty-aware selection.

\subsection{Sensitivity Analysis}

A $2^3$ factorial tests demand noise ($\sigma_d \in \{0.03, 0.15\}$), forced outage rates (all 1\% vs.\ realistic), and renewable variability (fixed mean vs.\ stochastic), restricted to the top-4 portfolios with 30 CRN replications per cell.

\subsection{Extensions}

Two additional analyses extend the core evaluation. First, a \textit{stochastic adequacy search} finds, for each top-4 mix, the minimum total capacity such that EUE $\leq 10$~MWh/year via bisection (50 replications per step). Second, a \textit{UCAP-normalized comparison} re-evaluates six portfolios (spanning 20--70\% firm) each sized to equal effective capacity (8,000~MW UCAP) using ELCC values from IESO planning conventions (Appendix~\ref{app:elcc}). This tests whether the equal-installed-MW design biases results toward firm-heavy portfolios. Finally, a \textit{post-processing VOLL sensitivity} recomputes $Y = F + C + \lambda U$ over $\lambda \in \{5{,}000, 10{,}000, 15{,}000\}$ using the same replication-level components from the extended CRN run.

%==========================================================================
% SECTION 5: RESULTS
%==========================================================================
\section{Results}\label{sec:results}

\subsection{Portfolio Selection}

The KN procedure selects P04 (70/20/10) after 77 replications (Figure~1). Fourteen of sixteen
portfolios are eliminated at the initial stage ($n_0 = 30$); only P01 (40/50/10) survives to
stage~77. Table~2 reports the extended evaluation over 200 CRN replications.

\begin{figure}[ht]
\centering
\includegraphics[width=0.97\textwidth]{figures/fig2_kn_elimination.png}
\caption{KN elimination timeline. Fourteen portfolios are eliminated at $n_0 = 30$; P01 survives to stage 77 before elimination. P04 (70/20/10) is selected as the best portfolio.}\label{fig:kn}
\end{figure}

\begin{table}[ht]
\centering
\caption{Top-4 portfolio evaluation (200 replications).}\label{tab:topk}
\begin{tabular}{lccccc}
\toprule
Portfolio & Mean $Y$ (\$M) & $F_i$ (\$M) & $\bar{C}$ (\$M) & $\lambda\bar{U}$ (\$M) & RHW \\
\midrule
P04 (70/20/10) & 4,306 & 1,694 & 1,835 & 777   & 0.38\% \\
P10 (70/10/20) & 5,060 & 1,732 & 1,874 & 1,454 & 0.40\% \\
P03 (60/30/10) & 8,272 & 1,660 & 1,846 & 4,766 & 0.83\% \\
P09 (60/20/20) & 10,785 & 1,698 & 1,882 & 7,205 & 0.64\% \\
\bottomrule
\end{tabular}
\end{table}
\begin{figure}[ht]
\centering
\includegraphics[width=0.95\textwidth]{figures/fig8_cost_decomposition.png}
\caption{Cost decomposition: fixed cost $F_i$ (blue), operating cost $C_{ij}$ (orange), and adequacy penalty $\lambda U_{ij}$ (red). P04 pays the highest fixed cost but the lowest penalty, yielding the lowest total.}\label{fig:decomp}
\end{figure}

The cost decomposition (Figure~2) reveals the mechanism driving this result. P04 incurs the
highest annualized fixed cost among the top-4 portfolios (\$1,694M), reflecting the capital
intensity of its 70\% firm allocation. However, its adequacy penalty is dramatically lower:
\$777M compared to \$4,766M for P03 and \$7,205M for P09. The adequacy penalty, not operating
cost, is the dominant source of performance differentiation. Operating costs vary by less than
3\% across the top-4 portfolios (\$1,835M to \$1,882M), confirming that the resource mix
primarily affects reliability rather than dispatch efficiency.

The physical mechanism is the interaction between low firm capacity and temporally correlated
renewable shortfalls. The NORTA AR(1) wind model ($\rho = 0.995$) produces sustained
multi-hour periods of low output that do not appear under iid sampling or deterministic
mean-input evaluation. When a portfolio holds insufficient firm generation to cover these
extended low-renewable events, unserved energy accumulates rapidly. P03 (60/30/10) has
1,200~MW less firm capacity than P04, and this deficit is exposed precisely during the
sustained wind droughts that the autocorrelated input model captures. The heavy right tails
visible in Figure~9 for P03 and P09 reflect these high-penalty replications.

All six pairwise differences among the top-4 portfolios are statistically significant under
Bonferroni correction ($\alpha_{\text{pair}} = 0.0083$), and all relative half-widths are
below 1\% (Table~2), confirming that the ranking is precisely estimated. The KN procedure
terminated at 77 replications, well below the 500-replication cap, reflecting the efficiency
of CRN in reducing pairwise difference variance (2$\times$ to 35$\times$, Figure~10).


\subsection{Deterministic Benchmark Gap}
Under deterministic assumptions (demand noise set to zero, all generators available, renewable
capacity factors at their annual means), the lowest-cost portfolio is P03 (60/30/10) at
\$3,503M. This differs from the stochastic recommendation (P04). The DBG is:
%
\begin{equation}\label{eq:dbg}
  \text{DBG} = \bar{Y}_{\text{P03}} - \bar{Y}_{\text{P04}} = \$3{,}968\text{M}, \quad 95\%\;\text{CI:}\; [\$3{,}905,\; \$4{,}031\text{M}].
\end{equation}
A planner relying on mean-input evaluation would select a portfolio costing approximately
\$4B more per year when uncertainty is included (Figure~3).

The source of the gap is informative. P03 saves \$34M in annualized fixed cost relative to P04
because its lower firm fraction requires less capital-intensive generation. Under deterministic
evaluation, this saving is not offset by any adequacy penalty, because mean renewable output is
sufficient to cover residual demand when no outages occur and no demand noise is present. The
deterministic evaluation effectively assumes the system always operates near its expected
conditions, rendering the 10 percentage-point difference in firm capacity invisible.

Under stochastic evaluation, the tail risk reappears. Demand noise ($\sigma_d = 0.09$)
occasionally pushes load above the level that P03's firm fleet can cover, and forced outages
(5--8\% FOR) reduce available firm capacity precisely when it is needed most. The
NORTA-modeled wind droughts compound this exposure: during multi-hour low-wind events,
P03's 3,600~MW of renewables contribute negligibly, and its 7,200~MW of firm generation
(reduced by outages) cannot bridge the gap. The \$4B DBG therefore quantifies the cost of
ignoring these correlated tail events in portfolio selection.

This finding directly addresses Research Question~1: the preferred portfolio changes when
uncertainty is included, and the cost consequence of ignoring this is substantial relative to
system-level annual expenditures.

\begin{figure}[ht]
\centering
\includegraphics[width=0.95\textwidth]{figures/fig6_dbg.png}
\caption{Deterministic Benchmark Gap. Under mean-input assumptions, P03 appears cheaper. Under stochastic evaluation with 200 paired CRN replications, P04 costs \$3,968M less per year.}\label{fig:dbg}
\end{figure}

\subsection{Sensitivity Analysis}

The $2^3$ factorial (Figure~\ref{fig:tornado}) identifies demand noise (\$2,883M main effect) and forced outages (\$2,838M) as the dominant uncertainty factors. Renewable variability contributes only \$117M, reflecting the concentrated Beta(6.73, 15.69) marginal calibrated from fleet-level CODERS data: because aggregate Ontario wind capacity factors cluster tightly around 0.30, switching between stochastic and fixed-mean renewable profiles has limited impact on the performance gap between the best and worst portfolios.

The demand$\times$outage interaction ($-$\$699M) is substantial and merits interpretation. A negative interaction means the combined effect of high demand noise and realistic forced outages is \emph{less} than the sum of their individual effects. The explanation is a saturation mechanism: when either factor alone is at its high level, the system is already under significant stress, and portfolios with low firm capacity are already incurring large adequacy penalties. Activating the second stressor increases unserved energy further, but the marginal impact is diminishing because the most vulnerable hours are already experiencing shortfalls. In practical terms, modeling demand uncertainty alone captures much of the adequacy risk that forced outages would additionally reveal, and vice versa.

The most important finding from the factorial is that P04 is the best portfolio in all eight configurations. This indicates the recommendation is robust to uncertainty assumptions within the tested range, directly addressing Research Question~2: while demand noise and forced outages materially affect the \emph{magnitude} of portfolio costs, they do not change the \emph{ranking}. The portfolio selection is therefore not sensitive to the specific calibration of these uncertainty factors, strengthening confidence in the stochastic recommendation.

\begin{figure}[ht]
\centering
\includegraphics[width=0.95\textwidth]{figures/fig4_tornado.png}
\caption{Factorial sensitivity analysis. Demand noise and forced outages dominate; renewable variability has minimal impact. P04 is the best portfolio in all eight configurations.}\label{fig:tornado}
\end{figure}

\subsection{VOLL Sensitivity}

To test robustness to the adequacy penalty coefficient, the top-4 portfolios were re-ranked for $\lambda \in \{5{,}000,10{,}000,15{,}000\}$ \$/MWh using the same replication-level $(F,C,U)$ outputs from the 200-replication CRN experiment. The selected portfolio remains unchanged (P04) across all tested values. This indicates the recommendation is not an artifact of a single VOLL calibration point, at least within a practically relevant range.

\subsection{Adequacy Search}

When total capacity is a decision variable, P04's mix achieves adequacy (EUE $\approx 0$) at 15.8~GW, while P03 requires 17.5~GW and P09 requires 18.2~GW. The cheapest adequate portfolio is P04 at \$3,763M (Appendix, Figure~\ref{fig:adequacy}).

\subsection{UCAP-Normalized Comparison}

To test whether equal-installed-MW comparisons bias toward high-firm portfolios, six mixes spanning 20--70\% firm are evaluated at equal UCAP (8,000~MW). Even after normalization, P04 remains cheapest (\$9,186M vs.\ \$10,167M for P03). The adequacy penalty dominates for low-firm mixes despite higher installed capacity, because ELCC is a static metric that does not capture sustained low-renewable events. During multi-hour wind droughts ($\rho = 0.995$), effective renewable UCAP drops to near zero, and additional installed capacity provides no incremental adequacy value (Appendix, Figure~\ref{fig:ucap}).

%==========================================================================
% SECTION 6: CONCLUSION
%==========================================================================
\section{Conclusion}\label{sec:conclusion}

This study demonstrates a simulation-based framework for comparing capacity portfolios under uncertainty, combining chronological Monte Carlo simulation with empirically calibrated inputs, CRN, and Kim--Nelson ranking and selection.

The central finding is that uncertainty changes the preferred portfolio. Deterministic analysis selects P03 (60/30/10); stochastic analysis selects P04 (70/20/10), which costs \$4B less per year when uncertainty is included (DBG = \$3,968M, $p < 0.05$). This is robust across all factorial configurations and persists under UCAP normalization.

Three methodological contributions merit emphasis. First, the NORTA AR(1) model preserves both the empirical marginal and temporal autocorrelation ($\rho = 0.995$) of Ontario wind data; iid sampling underestimates sustained low-output events. Second, CRN achieves 2--35$\times$ variance reduction, substantially accelerating KN elimination. Third, the adequacy search demonstrates stochastic root-finding for minimum reliable capacity, bridging evaluation with expansion planning.

\paragraph{Discussion.}
The results collectively address both research questions. Research Question~1 is answered by the deterministic benchmark gap: the preferred portfolio changes from P03 under deterministic evaluation to P04 under stochastic evaluation, with a cost consequence of \$3,968M per year. Research Question~2 is answered by the factorial analysis: demand noise and forced outages are the dominant uncertainty factors by main effect magnitude, but no tested factor combination changes the identity of the best portfolio. The recommendation is therefore robust to uncertainty calibration within the tested range.

These findings are consistent with the central insight of \citet{ParkBaldick2016}, who show that incorporating uncertainty into capacity planning can change the preferred resource mix. Their study demonstrates this for expansion timing decisions in a multi-stage stochastic program; the present study demonstrates the analogous result for portfolio ranking under a fixed-budget constraint. The methodological difference is consequential: whereas Park and Baldick jointly optimize build decisions and uncertainty realizations, this study takes candidate portfolios as given and applies formal statistical selection (KN with CRN) to identify the best performer. The shared conclusion is that deterministic planning can systematically undervalue firm capacity by failing to account for the tail risk introduced by correlated renewable shortfalls and generator outages.

\paragraph{Limitations.}
(1)~Demand noise is iid, which may overstate extreme event frequency.
(2)~Storage dispatch is rule-based, not co-optimized.
(3)~NORTA AR(1) assumes normal-space $\rho$ equals the target autocorrelation; the error is negligible at $\rho = 0.995$.
(4)~Discrete outage granularity (10 units per sub-type).

\paragraph{Future work.}
The natural extension is full stochastic capacity expansion, jointly optimizing the portfolio mix and total capacity subject to an adequacy constraint. The adequacy search solves a one-dimensional version; the multi-dimensional formulation connects directly to \citet{ParkBaldick2016}.

%==========================================================================
% REFERENCES
%==========================================================================
\bibliographystyle{plainnat}
\bibliography{references}

%==========================================================================
% APPENDIX
%==========================================================================
\appendix

\section{Consolidated Parameters}\label{app:params}

\begin{table}[H]
\centering
\caption{Simulation parameters.}\label{tab:params}
\small
\begin{tabular}{lll}
\toprule
Parameter & Value & Basis \\
\midrule
\multicolumn{3}{l}{\textit{System}} \\
Peak demand & 8,000 MW & Stylized \\
Reserve margin & 50\% & Accounts for renewable ELCC \\
Total capacity & 12,000 MW & Peak $\times$ 1.5 \\
\midrule
\multicolumn{3}{l}{\textit{Demand}} \\
Profile & 2-harmonic OLS & IESO 2024 ($R^2 = 0.59$) \\
Noise $\sigma_d$ & 0.09 & Residual std from fit \\
\midrule
\multicolumn{3}{l}{\textit{Renewables}} \\
Wind CF & Beta(6.73, 15.69) & MLE, CODERS Ontario \\
Wind $\rho$ & 0.9945 & Lag-1 ACF, CODERS \\
Solar cloud & Beta(1.77, 6.91) & MLE, CODERS Ontario \\
Solar $\rho$ & 0.928 & Lag-1 ACF, CODERS \\
Daylight hours & 07:00--21:00 & CODERS (UTC$-$5) \\
Wind/solar split & 60/40 & Stylized \\
\midrule
\multicolumn{3}{l}{\textit{Storage}} \\
Duration & 4 hours & Industry standard \\
Round-trip eff.\ & 85\% & NREL ATB \\
Discharge/charge & 70th/30th pctile & Stylized rules \\
\midrule
\multicolumn{3}{l}{\textit{Cost}} \\
VOLL ($\lambda$) & \$10,000/MWh & Mid-range \\
Nuclear capital & \$400/kW-yr & High overnight, 30-yr \\
Flex firm capital & \$100/kW-yr & Moderate \\
Mid-merit capital & \$70/kW-yr & Combined cycle \\
Peaker capital & \$70/kW-yr & Simple cycle \\
Wind capital & \$130/kW-yr & Onshore \\
Solar capital & \$100/kW-yr & Utility-scale \\
Storage capital & \$150/kW-yr & Battery, 15-yr life \\
\midrule
\multicolumn{3}{l}{\textit{KN procedure}} \\
$\alpha$ / $\delta$ / $n_0$ & 0.05 / \$50M / 30 & \\
\bottomrule
\end{tabular}
\end{table}


\section{Portfolio Definitions}\label{app:portfolios}

\begin{table}[H]
\centering
\caption{All 16 candidate portfolios (12,000 MW total).}\label{tab:allport}
\small
\begin{tabular}{clrrrr}
\toprule
ID & Label & Firm & Renew & Storage & $F_i$ (\$M) \\
\midrule
P00 & 30/60/10 & 3,600 & 7,200 & 1,200 & 1,557 \\
P01 & 40/50/10 & 4,800 & 6,000 & 1,200 & 1,591 \\
P02 & 50/40/10 & 6,000 & 4,800 & 1,200 & 1,625 \\
P03 & 60/30/10 & 7,200 & 3,600 & 1,200 & 1,660 \\
P04 & 70/20/10 & 8,400 & 2,400 & 1,200 & 1,694 \\
P05 & 20/60/20 & 2,400 & 7,200 & 2,400 & 1,561 \\
P06 & 30/50/20 & 3,600 & 6,000 & 2,400 & 1,595 \\
P07 & 40/40/20 & 4,800 & 4,800 & 2,400 & 1,630 \\
P08 & 50/30/20 & 6,000 & 3,600 & 2,400 & 1,664 \\
P09 & 60/20/20 & 7,200 & 2,400 & 2,400 & 1,698 \\
P10 & 70/10/20 & 8,400 & 1,200 & 2,400 & 1,732 \\
P11 & 20/50/30 & 2,400 & 6,000 & 3,600 & 1,600 \\
P12 & 30/40/30 & 3,600 & 4,800 & 3,600 & 1,634 \\
P13 & 40/30/30 & 4,800 & 3,600 & 3,600 & 1,668 \\
P14 & 50/20/30 & 6,000 & 2,400 & 3,600 & 1,702 \\
P15 & 60/10/30 & 7,200 & 1,200 & 3,600 & 1,736 \\
\bottomrule
\end{tabular}
\end{table}


\section{ELCC Values}\label{app:elcc}

\begin{table}[H]
\centering
\caption{ELCC values for UCAP normalization.}\label{tab:elcc}
\small
\begin{tabular}{lcc}
\toprule
Technology & ELCC & Basis \\
\midrule
Nuclear      & 94\% & High availability baseload \\
Flexible firm & 92\% & Gas/hydro, slight derates \\
Mid-merit gas & 95\% & CCGT \\
Gas peaker   & 92\% & Cycling stress \\
Wind         & 18\% & IESO capacity auction \\
Solar        & 35\% & Summer peak contribution \\
Storage (4h) & 50\% & Duration-limited \\
\bottomrule
\end{tabular}
\end{table}


\section{Input Fitting Diagnostics}\label{app:inputs}

\begin{figure}[H]
\centering
\includegraphics[width=\textwidth]{figures/input_demand_fit.png}
\caption{Demand profile OLS fit to IESO 2024 hourly Ontario demand. Left: observed vs.\ fitted. Right: residual distribution ($\sigma_d = 0.09$).}\label{fig:demand_fit}
\end{figure}

\begin{figure}[H]
\centering
\includegraphics[width=\textwidth]{figures/input_wind_fit.png}
\caption{Wind CF Beta(6.73, 15.69) MLE fit to CODERS Ontario. Left: empirical vs.\ fitted PDF. Right: Q--Q plot (KS $= 0.021$).}\label{fig:wind_fit}
\end{figure}

\begin{figure}[H]
\centering
\includegraphics[width=\textwidth]{figures/input_solar_fit.png}
\caption{Solar CF: diurnal envelope $\times$ Beta(1.77, 6.91), calibrated from CODERS Ontario.}\label{fig:solar_fit}
\end{figure}


\section{Supplementary Results}\label{app:results}

\begin{figure}[H]
\centering
\includegraphics[width=0.95\textwidth]{figures/fig1_mean_y_bar.png}
\caption{Mean $Y$ with 95\% CIs for top-4 portfolios. The firm-fraction gradient drives cost: portfolios with less firm capacity incur large adequacy penalties.}\label{fig:means}
\end{figure}

\begin{figure}[H]
\centering
\includegraphics[width=0.95\textwidth]{figures/fig3_y_distributions.png}
\caption{Distribution of $Y_{ij}$ for top-4 portfolios. P03 and P09 exhibit heavy right tails from high-penalty replications during sustained low-renewable periods.}\label{fig:box}
\end{figure}

\begin{figure}[H]
\centering
\includegraphics[width=0.95\textwidth]{figures/fig7_crn_efficiency.png}
\caption{CRN variance reduction: pairwise difference variance under independent vs.\ CRN sampling. Efficiency ranges from 2$\times$ (P04 vs.\ P09) to 35$\times$ (P03 vs.\ P09).}\label{fig:crn}
\end{figure}

\begin{figure}[H]
\centering
\includegraphics[width=0.95\textwidth]{figures/fig10_adequacy_search.png}
\caption{Stochastic adequacy search. Left: minimum capacity for EUE $\approx 0$. Right: annualized fixed cost at adequate capacity. P04 achieves adequacy at the lowest capacity and cost.}\label{fig:adequacy}
\end{figure}

\begin{figure}[H]
\centering
\includegraphics[width=0.95\textwidth]{figures/fig12_ucap_comparison.png}
\caption{UCAP-normalized comparison. Even at equal effective capacity (8~GW UCAP), P04 remains cheapest. Low-firm portfolios incur large adequacy penalties because ELCC does not capture sustained low-renewable events.}\label{fig:ucap}
\end{figure}


\section{Code Architecture}\label{app:code}

The simulation is implemented in Python 3.11 (NumPy, SciPy). Twelve source files total approximately 2,500 lines. The full pipeline runs in approximately 5 minutes on a consumer laptop. All results are reproducible via \texttt{python src/main.py}.

\begin{table}[H]
\centering
\caption{Code modules.}\label{tab:code}
\small
\begin{tabular}{ll}
\toprule
Module & Responsibility \\
\midrule
\texttt{config.py} & Parameters, portfolio generation, ELCC \\
\texttt{crn.py} & CRN via SeedSequence, NORTA AR(1) \\
\texttt{dispatch.py} & Hourly merit-order dispatch \\
\texttt{kn.py} & Kim--Nelson sequential elimination \\
\texttt{experiment.py} & DBG, factorial, Bonferroni CIs, VOLL sensitivity \\
\texttt{enhancements.py} & CRN efficiency, decomposition, RHW \\
\texttt{input\_fitting.py} & Parametric fitting to CODERS/IESO \\
\texttt{adequacy.py} & Stochastic adequacy search \\
\texttt{ucap\_comparison.py} & UCAP-normalized comparison \\
\texttt{plots.py} & Figure generation \\
\texttt{main.py} & Pipeline orchestration (10 stages) \\
\bottomrule
\end{tabular}
\end{table}


\section{Generative AI Disclosure}\label{app:ai}

Generative AI (Claude, Anthropic) was used for assistance with code development, debugging, statistical methodology review, and drafting support. All modeling decisions, parameter choices, experimental design, interpretation of results, and final writing are the author's own.

\end{document}
